\documentclass{pgnotes}

\title{XML}
\newcommand{\entityref}[1]{\&#1;}

\tikzstyle{cpt} = [fill=black, text=white, draw=none, rounded corners, rectangle, minimum width=2.5cm, minimum height=1.5cm]
\tikzstyle{ml} = [draw=black, ultra thick, rounded corners, rectangle, minimum width=2.75cm, minimum height=1.5cm] 
\tikzstyle{subset} = [draw=black, ultra thick, ->, dashed] 
\tikzstyle{application} = [draw=black, ultra thick, ->] 
\tikzstyle{bizr} = [anchor=south, text=white] 


\begin{document}

\maketitle

% \section*{Aims}

% \begin{enumerate}
%   \item Explain the essential structures and rules of an XML document.
%   \item Distinguish well-formed XML from non-well-formed XML. Use tools to diagnose non-well-formed XML.
%   \item Construct vocabularies to capture structured information using XML.
% \end{enumerate}


\section{eXtensible Markup Language}

Extensible Markup Language (XML) is a \textbf{markup} language that defines a set of rules for encoding documents in a format that is both human-readable and machine-readable.
It's a subset of SGML, of which HTML is also an \textit{application}. 
The tag names are up to you. 

\subsection{SGML}

Standard Generalised Markup Language (SGML) originated as a means to store and transfer US government and military documents in machine-readable form, over a timespan of several decades.
SGML uses \textbf{tags} to mark up a document.
One of the best known examples of SGML is HTML.

\inputminted{html}{htmlpage.html}

For many applications, SGML was too complicated, and so XML was designed.
Apart from HTML and Document Type Declarations, you are unlikely to encounter SGML in its non-XML form very frequently.
XML documents themselves are valid SGML, but have additional restrictions on their layout.
(SGML documents are not valid XML.)
 

\subsection{Sample XML document}

\inputminted{xml}{intro.xml}

\subsection{Design Goals}

\begin{enumerate}
\item XML shall be straightforwardly usable over the Internet.
\item XML shall support a wide variety of applications.
\item XML shall be compatible with SGML.
\item It shall be easy to write programs which process XML documents.
\item The number of optional features in XML is to be kept to the absolute minimum, ideally zero.
\item XML documents should be human-legible and reasonably clear.
\item The XML design should be prepared quickly.
\item The design of XML shall be formal and concise.
\item XML documents shall be easy to create.
\item Terseness in XML markup is of minimal importance.
\end{enumerate}

 

\subsection{Benefits of XML}

\begin{enumerate}
\item You define the \textbf{vocabulary}, so it provides flexibility.
\item XML is \textbf{simple} and used widely.
\item Machine and \textbf{human readable}.
\item \textbf{Lots of APIs} and parsers to handle XML. 
\item XML Data \textbf{can be checked} for correctness.
\end{enumerate}

\subsection{Issues with XML}

\begin{enumerate}
\item XML documents require more storage than fixed length or comma delimited files.
  \begin{itemize}
  \item But, they compress well!
  \end{itemize}
\item Can help but not does not solve the politics of competing formats (multiple vocabs).
\item Definitions can be overly nested and made too complex.
\end{enumerate}
 
\section{Structure of an XML doc}

XML has a Directed Acyclic Graph tree structure, built from \textbf{nested} tree of \textbf{Elements} delimited by \textbf{tags}.
Elements are \textbf{case-sensitive}.

\inputminted{xml}{closed.xml}

\subsection{Nesting rules} 

\begin{itemize}
\item Elements nested inside a \textit{parent} element are called \textit{child} elements. 
\item Elements must be nested correctly. Child elements must be enclosed within their parent elements.
\item All elements must be nested within a single document or root element.
\item There can be only one \textbf{root} element in a document.
\item Indentation, while helpful, is just visual.  
\end{itemize}


\subsection{XML declaration}

The XML declaration is always the first line of code in an XML document, and tells the \textbf{processor} that what follows is XML.
It can also provide information about how the parser should interpret the document.

For now, the version is always 1.0.
This attribute allows the program (or person) reading the document to read older versions if specifications change. 


\subsection{XML comments}

Comments may appear anywhere after the declaration.
The syntax for comments is:

\inputminted{xml}{comment.xml}

Two or more dashes should not appear one after the other in a comment.


\subsection{Free-format}

XML documents are free-format.
For example, the document:

\inputminted{xml}{free-format-condensed.xml}
  
is the exact same as:
  
\inputminted{xml}{free-format-expanded.xml}


\section{Attributes}

An attribute is a feature or characteristic of an element. Attributes are text strings and must be placed in single or double quotes. 
Use that which makes most sense in deciding whether to store the main data in an element or attribute.
Often an element is a better choice.

\inputminted{xml}{attribute.xml}


\subsection{Mixing elements and attributes}

\inputminted{xml}{mixing.xml}



\section{Well-formed XML}

When XML text satisfies the syntactic rules as laid out in the XML specification, we say that it is well-formed \citep{goldfarb:2003:the-xml-handbook}.

\subsection{Criteria}

\begin{enumerate}
\item All XML elements must have a closing tag, except empty elements. 
\item Tags are case sensitive.
\item All elements must be properly nested.
\item Documents must have one single root element.
\item Attribute values must be quoted.
\item Certain characters must be replaced with \textit{Entity References}. 
\end{enumerate}


\subsection{Entity references}

Entity references replace certain characters as detailed in \autoref{tab:entity-reference-xml}.

\begin{table}[htbp]
  \centering
  \begin{tabular}{l l l}
    \toprule
    \textbf{Normal symbol} & \textbf{Meaning} & \textbf{Replace with} \\
    \midrule
    $<$ & less than & \entityref{lt} \\
    $>$ & greater than & \entityref{gt} \\
    \& & ampersand & \entityref{amp} \\
    ' & apostrophe & \entityref{apos} \\
    '' & quotation mark & \entityref{quot} \\
    \bottomrule
  \end{tabular}
  \caption{Entity References}
  \label{tab:entity-reference-xml}
\end{table}


\subsection{Non well-formed XML}

Non well-formed XML will cause errors in applications and parsers that expect strict XML syntax.

\subsection{Line termination}

By convention, different OSes use different delimiters to signal the end of a line in a text file, \autoref{tab:newline}.
XML specifies that new lines are \textbf{always} stored as LF, regardless of the operating system. 

\begin{table}[htbp]
  \centering
  \begin{tabular}{l l}
    \toprule
    \textbf{Operating system} & \textbf{Delimiters} \\
    \midrule
    Windows & CR LF \\
    Unix, Mac OS X & LF \\
    Mac (System 9 and below) & CR \\ 
    \bottomrule
  \end{tabular}
  \caption{Newline characters}
  \label{tab:newline}
\end{table}

\subsection{Checking for well-formed XML}

The \texttt{xmllint} utility can check XML well-formedness.

\begin{minted}{bash}
# consult manpage
man xmllint

# basic usage
xmllint filename.xml 
\end{minted}
 
  








\section{Vocabulary}

Although the structure and syntax of XML documents is rigorously specified, you are free to choose a suitable vocabulary of element and attribute names for the application of interest.
Some common vocabularies are listed in \autoref{tab:common-vocabularies}. 

\begin{table}[htbp]
  \centering
  \begin{tabularx}{1.0\linewidth}{l X}
    \toprule
    \textbf{Vocabulary} & \textbf{Description} \\
    \midrule
    XHTML & a variant of HTML that complies with XML restrictions \\
    SOAP & message format for web services \\
    OpenOffice XML & Microsoft's current office word-processing format \\
    Opendoc & standard formats for (non-Microsoft Office) office documents \\
    Docbook XML & markup for writing (primarily) technical documentation  \\
    \bottomrule
  \end{tabularx}
  \caption{Common standardised XML vocabularies}
  \label{tab:common-vocabularies}
\end{table}

In addition, there are pre-existing vocabularies for many specialised applications.
It's always good to use a pre-existing vocabulary if possible. 

\subsection{Developing an XML vocabulary}

The process is:

\begin{enumerate}
\item List out the possible fields (tag names).
\item Draw a tree diagram.
\item Markup some sample XML data.
\item Factor in potential business rules and add them to our tree diagram. 
\item Optionally formalise the vocabulary in a machine-readable form.
\end{enumerate}

\subsection{Worked example}

As an example, we will develop a vocabulary that would let you specify address data for a person.

\subsubsection{Possible fields}
Start with:
\begin{itemize}
\item firstname
\item secondname
\item address1
\item address2
\item city
\item country
\end{itemize}

Each of these can be specified as tags like
\begin{minted}{xml}
<firstname></firstname>
\end{minted}

But: a person's name is not really part of an address.
What happens if a person has multiple addresses?
So we really need a nested structure:
\begin{minted}{xml}
<person>
  <address>
  </address>
</person>
\end{minted}


\subsubsection{Diagram}

The nested structure of our document is then laid out as a \textbf{tree diagram}, \autoref{fig:tree-diagram}. 

\begin{figure}[htbp]
  \centering
  \begin{tikzpicture}
  \node (person) [cpt] {person};
  \node (firstname) [cpt, below left=of person] {firstname};
  \draw (person) to (firstname);
  \node (secondname) [cpt, below=of person] {secondname};
  \draw (person) to (secondname);
  \node (address) [cpt, below right=of person] {address};
  \draw (person) to (address); 
  \node (street) [cpt, below left=of address] {street};
  \draw (address) to (street); 
  \node (city) [cpt, below=of address] {city};
  \draw (address) to (city); 
  \node (country) [cpt, below right=of address] {country}; 
  \draw (address) to (country); 
  \end{tikzpicture}
  \caption{Tree diagram for address example}
  \label{fig:tree-diagram}
\end{figure}

\subsubsection{Sample data}

\inputminted{xml}{sampledata.xml}

\subsubsection{Rules}

We then decide on a number of business rules to determine if an element is optional, and if there are any constraints on its \textbf{multiplicity}: 

\begin{itemize}
\item Person can have multiple first names.
\item Person can have only one second names.
\item Person can have multiple addresses.
\item An address can have a few streets (lets say between 1 and 3).
\end{itemize}

These rules are then added to the tree diagram, \autoref{fig:tree-diagram-with-rules}. 

\begin{figure}[htbp]
  \centering
  \begin{tikzpicture}
  \node (person) [cpt] {person};
  \node (firstname) [cpt, below left=of person, fill=DarkBlue] {firstname};
  \node [bizr] at (firstname.south) {min=1 max=*};  
  \draw (person) to (firstname);
  \node (secondname) [cpt, below=of person] {secondname};
  \draw (person) to (secondname);
  \node (address) [cpt, below right=of person, fill=DarkBlue] {address};
  \node [bizr] at (address.south) {min=1 max=*};  
  \draw (person) to (address); 
  \node (street) [cpt, below left=of address, fill=DarkGreen] {street};
  \node [bizr] at (street.south) {min=1 max=3};  
  \draw (address) to (street); 
  \node (city) [cpt, below=of address] {city};
  \draw (address) to (city); 
  \node (country) [cpt, below right=of address] {country}; 
  \draw (address) to (country); 
  \end{tikzpicture}
  \caption{Tree diagram with business rules for address example}
  \label{fig:tree-diagram-with-rules}
\end{figure}
 


\bibliographystyle{plainnat}
\bibliography{bibliography}

\end{document}